\section{Matrices and Linear Maps}

Let \(V, W\) be linear maps. Define the set
\[
    \linearmap\pqty{V, W} = \set{\functiondomain{\alpha}{V}{W}}{\alpha \text{ is a linear map}}
\]
equipped with element-wise addition and scalar multiplication
\[
    \pqty{\alpha + \beta} \pqty{\vb{v}} = \alpha\pqty{\vb{v}} + \beta\pqty{\vb{v}} \qand \pqty{\lambda \alpha} \pqty{\vb{v}} = \lambda \alpha \pqty{\vb{v}}.
\]

It is straightforward to check that \(\linearmap\pqty{V, W}\) is an \(\FF\)-vector space.

Suppose onwards that \(V\) and \(W\) are finite-dimensional, with \(\dim \pqty{V} = n\) and \(\dim \pqty{W} = m\).

\begin{theorem}[Linear Maps and Matrices]
    \label{thm:linear-maps-matrices}%
    We have
    \[
        \linearmap\pqty{V, W} \isomorphic \matrixset_{m \times n} \pqty{\FF}.
    \]

    In particular, we have \(\dim\pqty{\linearmap\pqty{V, W}} = mn\).
\end{theorem}

\begin{notation}
    For \(A \in \matrixset_{m \times n} \pqty{\FF}\), we denote the entries of \(A\) as \(A = \pqty{a_{i, j}}\) where \(a_{i, j} \in \FF\) and \(1 \leq i \leq m\), \(1 \leq j \leq n\). Furthermore, let \(\matrixset_{m \times n}\pqty{\FF}\) has a standard basis
    \[
        \set{E_{r, s}}{1 \leq r \leq m, 1 \leq s \leq n}
    \]
    where \(E_{r, s} = \pqty{\kronecker_{r, i}, \kronecker_{s, j}}\)
\end{notation}

Since \(V\) and \(W\) are finite dimensional, let \(\basis = \Bqty{\vb{v}_1, \ldots, \vb{v}_n}\) and \(\mathcal{C} = \Bqty{\vb{w}_1, \ldots, \vb{w}_m}\). Then given a linear map \(\alpha \in \linearmap\pqty{V, W}\), we may apply \(\alpha\) to the basis vectors \(\vb{v}_j\) to get linear combinations of \(\vb{w}_i\)
\begin{align*}
    \alpha\pqty{\vb{v}_1} &= a_{1, 1} \vb{w}_1 + a_{2, 1} \vb{w}_2 + \cdots + a_{m, 1} \vb{w}_m, \\
    \alpha\pqty{\vb{v}_2} &= a_{1, 2} \vb{w}_1 + a_{2, 2} \vb{w}_2 + \cdots + a_{m, 2} \vb{w}_m, \\
    \vdots &= \vdots \\
    \alpha\pqty{\vb{v}_n} &= a_{1, n} \vb{w}_1 + a_{2, n} \vb{w}_2 + \cdots + a_{m, n} \vb{w}_m.
\end{align*}

\begin{definition}[Matrix of Linear Map]
    The \term{matrix of the linear map} \(\alpha\) with respect to \(\basis\) and \(\mathcal{C}\) is given by
    \[
        \bqty{\alpha}_{\mathcal{C}}^{\basis} = \pqty{a_{i, j}} \in \matrixset_{m \times n} \pqty{\FF}
    \]
\end{definition}

The purpose of defining the matrix this way around is for consistent matrix multiplication with coordinate vectors. In particular, for all \(\vb{v} \in V\), we have
\[
    \bqty{\alpha}_{\mathcal{C}}^{\basis} \bqty{\vb{v}}_{\basis} = \bqty{\alpha\pqty{\vb{v}}}_{\mathcal{C}}.
\]

Furthermore, the matrix defined above is in fact the \emph{only} matrix \(A \in \matrixset_{m \times n} \pqty{\FF}\) that gives
\[
    A \bqty{\vb{v}}_{\basis} = \bqty{\alpha\pqty{\vb{v}}}_{\mathcal{C}}
\]
for all \(\vb{V} \in V\).

\begin{proof}
    \begin{enumerate}[label=(\roman*)]
        \item \textbf{Compatibility with Coordinate Vectors.} For any vector \(\vb{v} = \sum_{j = 1}^{n} \lambda_j \vb{v}_j\), we have
        \begin{align*}
            \alpha\pqty{\vb{v}} &= \sum_{j = 1}^{n} \lambda_j \alpha\pqty{\vb{v}_j} \\
            &= \sum_{j = 1}^{n} \lambda_j \sum_{i = 1}^{m} a_{ij} \vb{w}_i \\
            &= \sum_{i = 1}^{m} \pqty{\sum_{j = 1}^{n} a_{i, j} \lambda_j} w_i.
        \end{align*}
    
        Then note
        \[  
            \bqty{\alpha\pqty{\vb{v}}}_{\mathcal{C}} = \pqty{\sum_{j = 1}^{n} a_{ij} \lambda_j}_i = \pqty{a_{i, j}}_{i, j} \pqty{\lambda_j}_j = \bqty{\alpha}_{\mathcal{C}}^{\basis} \bqty{\vb{v}}_{\basis}
        \]
        is indeed the standard multiplication of the matrix and the basis coordinate vector.
    
        \item \textbf{Uniqueness with Coordinate Vectors}. By definition, \(\bqty{\vb{v_i}}_{\basis} = \vb{e}_i\) gives the standard basis vectors, so for \(A \in \matrixset_{m \times n} \pqty{\FF}\), we have
        \[
            A \bqty{\vb{v}_i}_{\mathcal{B}} = A \vb{e}_i = i\text{-th column of }A.
        \]
    
        But if another matrix \(A\) maps the corresponding coordinate vectors to the coordinate vectors, then
        \begin{align*}
            A \bqty{\vb{v}_i}_{\basis} &= \bqty{\alpha \pqty{\vb{v}_\mathcal{C}}}_{\mathcal{C}} \\
            &= \bqty{\alpha}_{\mathcal{C}}^{\basis} \bqty{\vb{v}_i}_{\basis} \\
            &= \bqty{\alpha}_{\mathcal{C}}^{\basis} \vb{e}_i \\
            &= i\text{-th colunn of } \bqty{\alpha}_{\mathcal{C}}^{\basis}
        \end{align*}
        and this holds for all \(i\), so \(A = \bqty{\alpha}_{\mathcal{C}}^{\basis}\).
    
        \item \textbf{Linearity.} For \(\alpha, \beta \in \linearmap\pqty{V, W}\) and \(\lambda \in \FF\), we can compare the matrices \(\bqty{\alpha}_{\mathcal{C}}^{\basis} + \lambda \bqty{\alpha}_{\mathcal{C}}^{\basis}\) with \(\bqty{\alpha + \lambda \beta}_{\mathcal{C}}^{\basis}\). For \(\vb{v} \in V\), we have
        \begin{align*}
            \pqty{\bqty{\alpha}_{\mathcal{C}}^{\basis} + \lambda \bqty{\beta}_{\mathcal{C}}^{\basis}} \bqty{\vb{v}}_{\basis} &= \bqty{\alpha \pqty{\vb{v}}}_{\mathcal{C}} + \lambda \bqty{\beta\pqty{\vb{v}}}_{\mathcal{C}} \\
            &= \bqty{\alpha \pqty{\vb{v}} + \lambda \beta \pqty{\vb{v}}}_{\mathcal{C}} \\
            &= \bqty{\pqty{\alpha + \lambda \beta}}_{\mathcal{C}}^{\basis} \bqty{\vb{v}}_{\basis}.
        \end{align*}
    
        Hence, by uniqueness in (ii), we have \(\bqty{\alpha + \lambda \beta}_{\mathcal{C}}^{\basis} = \bqty{\alpha}_{\mathcal{C}}^{\basis} + \lambda \bqty{\beta}_{\mathcal{C}}^{\basis}\).

        \item \textbf{Injectivity.} If \(\alpha, \beta \in \linearmap\pqty{V, W}\), with \(\bqty{\alpha}_{\mathcal{C}}^{\basis} = \bqty{\beta}_{\mathcal{C}}^{\basis}\), then we have
        \[  
            \alpha\pqty{\vb{v}_i} = \beta\pqty{\vb{v}_i}
        \]
        and by \autoref{prop:linear-map-is-determined-by-basis} (Uniqueness of Linear Extensions), we have \(\alpha = \beta\).

        \item \textbf{Surjectivity.} Given \(A \in \matrixset_{m \times n} \pqty{\FF}\), let \(\alpha in \linearmap\pqty{V, W}\) be the linear extension of
        \[
            \alpha\pqty{\vb{v}_j} = \sum_{i = 1}^{m} a_{i, j} \vb{w}_u
        \]
        with \(A = \pqty{a_{i, j}}\). Then \(\bqty{\alpha}_{\mathcal{C}}^{\basis} = A\).
    \end{enumerate}
\end{proof}

Using the \term{commutative diagram} below we can see what is going on here, with all the vector spaces.

\begin{figure}[ht!]
    \centering
    \begin{tikzcd}
        \vb{v} \arrow[rrr, mapsto] \arrow[ddd, mapsto] & & & \alpha\pqty{\vb{v}} \arrow[ddd, mapsto] \\
        & V \arrow[r, "\alpha"] \arrow[d, "F_{\basis}"] & W \arrow[d, "F_{\mathcal{C}}"] \\
        & \FF^n \arrow[r, "\bqty{\alpha}_{\mathcal{C}}^{\basis}"] & \FF^m \\
        \bqty{\vb{v}_{\basis}} \arrow[rrr, mapsto] & & & 
        \substack{\displaystyle \bqty{\alpha}_{\mathcal{C}}^{\basis} \bqty{\vb{v}}_{\basis} \\
        \displaystyle || \\
        \displaystyle \bqty{\alpha\pqty{\vb{v}}}_{\mathcal{C}}}
    \end{tikzcd}
    \caption{Commutative Diagram for Matrix and Linear Maps}
    \label{fig:commutative-diagram-matrix-linear-maps}%
\end{figure}

\begin{proposition}[Matrix of Composition]
    \label{prop:composition-matrix}%
    Let \(V, W, X\) be finite-dimensional vector spaces with bases \(\basis, \mathcal{C}, \mathcal{D}\) respectively, with \(\alpha \in \linearmap\pqty{V, W}\) and \(\beta \in \linearmap\pqty{W, X}\). Then
    \[
        \bqty{\beta \circ \alpha}_{\mathcal{D}}^{\basis} = \bqty{\beta}_{\mathcal{D}}^{\mathcal{C}} \bqty{\alpha}_{\mathcal{C}}^{\basis}.
    \]
\end{proposition}

\begin{proof}
    By \autoref{thm:linear-maps-matrices}, the matrix \(A = \bqty{\beta \circ \alpha}_{\mathcal{C}}^{\basis}\) is uniquely determined by
    \[  
        A \bqty{\vb{v}}_{\basis} = \bqty{\beta\pqty{\alpha\pqty{\vb{v}}}}_{\mathcal{D}}
    \]
    for all \(\vb{v} \in V\). But notice that
    \begin{align*}
        \pqty{\bqty{\beta}_{\mathcal{D}}^{\mathcal{C}} \bqty{\alpha}_{\mathcal{C}}^{\basis}} \bqty{\vb{v}}_{\basis} &= \bqty{\beta}_{\mathcal{D}}^{\mathcal{C}} \bqty{\alpha\pqty{\vb{v}}}_{\mathcal{C}} \\
        &= \bqty{\beta\pqty{\alpha\pqty{\vb{v}}}}_{\mathcal{D}}
    \end{align*}
    for all \(\vb{v} \in V\). Then we have \(A = \bqty{\beta}_{\basis}^{\mathcal{C}} \bqty{\alpha}_{\mathcal{C}}^{\basis}\).
\end{proof}

It is relatively obvious that for any basis \(\basis\) of \(V\), we have
\[
    \bqty{\identity_V}_{\basis}^{\basis} = \identityM_n.
\]

\begin{definition}[Base-Change Matrix]
    Given bases \(\basis, \basis'\) of \(V\), the \term{base-change matrix} from \(\basis\) to \(\basis'\) is given by
    \[
        P = \bqty{\identity_V}_{\basis'}^{\basis} = \matrixset_{m \times n} \pqty{\FF}.
    \]

    Equivalently, for \(\basis = \Bqty{\vb{v}_1, \vb{v}_2, \ldots, \vb{v}_n}\) and \(\basis' = \Bqty{\vb{v}_1', \vb{v}_2', \ldots, \vb{v}_n'}\), then \(P = \pqty{p_{i, j}}\), where
    \[
        \vb{v}_j = \sum_{i = 1}^{n} p_{i, j} \vb{v}_i'
    \]
    so the \(j\)-th column of \(P\) is \(\bqty{\vb{v}_j}_{\basis'}\).
\end{definition}

Recall that a matrix \(A \in \matrixset_{n \times n} \pqty{\FF}\) is \term{invertible} if there exists \(A\inverse \in \matrixset_{m \times n} \pqty{\FF}\) such that \(A A\inverse = A\inverse A = I_n\), and we call \(A\inverse\) an \term{inverse} for \(A\). The \term{group} of invertible \(n \times n\) matrices over \(\FF\) is denoted \(\GL_n \pqty{\FF}\).

\begin{proposition}[Change of Basis Matrices]
    We have the following.

    \begin{enumerate}[label=(\roman*)]
        \item The matrix \(\bqty{\identity_V}_{\basis'}^{\basis}\) is invertible, with inverse \(\bqty{\identity_V}_{\basis}^{\basis'}\).
        \item \term{Change of basis formula.} For \(\alpha \in \linearmap\pqty{V, W}\) and bases \(\basis, \basis'\) for \(V\), \(\mathcal{C}, \mathcal{C}'\) for \(W\), we have
        \[
            \bqty{\alpha}_{\mathcal{C}'}^{\basis'} = \bqty{\identity_W}_{\mathcal{C}'}^{\mathcal{C}} \bqty{\alpha}_{\mathcal{C}}^{\basis} \bqty{\identity_V}_{\basis}^{\basis'}.
        \]
    \end{enumerate}
\end{proposition}

\begin{proof}
    Both are immediate from \autoref{prop:composition-matrix}.
\end{proof}

Note that the converse is true: any invertible matrix is a base-change matrix.